\subsection{Equal and Proportional Rows}

If the two rows are equal, exchanging them leaves the matrix unchanged, but the row-exchange theorem changes the sign of the determinant. Hence
\[
\det A=-\det A,
\]
so $\det A=0$.

\begin{proofbox}[Direct check for proportional rows]
If
\[
A=\begin{pmatrix}a&b\\ka&kb\end{pmatrix},
\]
then
\begin{align*}
\det A
&=a(kb)-b(ka)\\
&=kab-kab\\
&=0.
\end{align*}
The six-term formula produces the same cancellation for $3\times3$ matrices with proportional rows.
\end{proofbox}

\begin{theorembox}[Dependent rows force zero]
For the small determinants already defined, equal or proportional rows give determinant zero.
\end{theorembox}

\begin{interpretationbox}[A relation removes freedom]
When one row repeats or rescales another, the rows do not carry independent information. The determinant records this loss by vanishing.
\end{interpretationbox}
